% 合成输入；各场景独立，不代表真实论文结果或学校规范。

% 场景 1
% 请求：请把这句话润色得准确、清楚，给出可替换的句子。
% 材料：没有其他结果记录。
该方法可能有助于提升预测性能。

% 场景 2
% 请求：请按 AXES 检查并改写这段结果分析。
% 材料：仅有本段数值，没有比较基线或组件试验。
注意力模型在测试集 T 上的准确率为95\%。这一提升源于模型捕获长程依赖的能力。结果见\ref{tab:single}。

% 场景 3
% 请求：请检查这段消融分析，给出局部改写建议。
% 材料：同一数据划分，完整模型训练100轮，去除模块A的模型训练50轮；其余训练条件未记录。
消融实验中，完整模型的准确率为95\%，去除模块A后为91\%。这证明模块A通过抑制噪声带来性能提升。

% 场景 4
% 请求：请检查结论强度，并给出局部改写建议。
% 材料：同一框架、数据划分、训练预算、优化器与种子集合，仅移除模块A。配对重复记录中，完整模型平均准确率95\%，移除后91\%；未观测中间噪声变量，也没有区分噪声抑制与其他机制的试验。
移除模块A后平均准确率由95\%降至91\%，说明模块A的噪声抑制机制导致了该差异。

% 场景 5
% 请求：请审阅以下结论能否保留，不需要重写已成立的句子。
% 材料：在协议P和测试块T内，以随机受控干预启用或关闭噪声通路N；数据、训练预算及其余模块固定，配对重复结果方向一致。中介测量和恢复通路试验排除了预先列出的容量变化及训练时长解释，支持噪声通路N对该协议下准确率差异的因果作用。
在协议P和测试块T内，准确率差异由噪声通路N的启用带来。通路定义为$z=x+\epsilon$，证据见\ref{tab:intervention}。

% 场景 6
% 请求：请润色这句贡献陈述。
% 材料：作者没有检索相关文献，没有优先权证据；X、Y为给定技术对象。
本文首次将X用于Y。

% 场景 7
% 请求：请检查摘要和章节安排是否适合当前博士论文，只提出诊断，不扩写正文。
% 材料：作者提供的本校要求为摘要1000—1500字、按研究工作编号，可定性陈述理论成果。当前摘要实测1100字，有对象背景与编号的两个工作段。第一项工作给出在假设H下的收敛证明，第二项为在CPU上运行的工程系统，只有离线日志回放证据。全文没有GPU训练或模型超参数搜索，不打算增加实验。
摘要工作段分别陈述假设H下的收敛结论与CPU工程系统的离线日志回放。第3章给出定理及证明，第4章按运行约束、接口机制与回放证据组织。收敛条件为$0<\eta<1$，理论定义见\cite{definition}。

% 场景 8
% 请求：请作一次逆向提纲审阅，定位段主题、章目标、证据与处置；只检查，不改正文。
% 材料：第2章目标是比较时序重建方法的输入假设并推导本文问题。当前窗口为2.1.2；prev.tail和next.head仅供理解衔接。下面给出源位置。
% prev.tail，chapters/review.tex:18
下节比较上述方法对输入完整性的要求。
% current，chapters/review.tex:24
现有时序重建方法对输入完整性的假设不同。固定采样方法要求等间隔观测，掩码建模方法显式接收缺失位置\cite{regular,masked,survey}。两类输入条件划定了本章比较的范围。
% current，chapters/review.tex:29
方法B在数据集D上的准确率为95\%，证明其在所有缺失条件下都能重建真实动态。
% current，chapters/review.tex:33
本文界面提供深色配色和菜单折叠选项，见\ref{fig:ui}。
% next.head，chapters/review.tex:38
本节据此分析不规则观测下的输入定义。
